%% ============================================================
%%  CHAPITRE 1 — Introduction à l'orchestration d'agents
%%               et au Model Context Protocol
%% ============================================================
\chapter{Introduction à l'orchestration d'agents et au Model Context Protocol}

Dans la continuité de la vue d'ensemble présentée en introduction, ce chapitre établit
les fondements théoriques et techniques nécessaires à la couche d'IA proposée. Nous
commençons par clarifier ce que signifie l'orchestration d'agents dans les systèmes
modernes d'IA générative, explorons les principaux schémas d'architecture actuellement
utilisés, puis nous concentrons sur le Model Context Protocol, l'innovation clé qui
rend praticable et passable à l'échelle l'intégration fiable avec des services backend existants.
La discussion s'approfondit ensuite sur des mécanismes d'orchestration concrets, en
particulier le schéma de superviseur LangGraph, et sur une vue détaillée de la
conception protocolaire et des propriétés de sécurité de MCP. Enfin, nous introduisons
le paradigme émergent de RAG agentique et passons en revue l'état de l'art des
applications en éducation afin de montrer comment ces concepts alimentent directement
notre implémentation SprintUp.

%% ────────────────────────────────────────────────────────────
\section{Comprendre l'orchestration d'agents}

Fondamentalement, un agent IA est plus qu'un simple chatbot. C'est un système autonome
qui perçoit son environnement (via les requêtes utilisateur ou les données récupérées),
raisonne sur des objectifs, sélectionne les outils ou actions appropriés et les exécute
en séquence ou en parallèle jusqu'à l'atteinte de l'objectif. Lorsque plusieurs agents
ou outils doivent collaborer, ce qui est inévitable pour traiter des demandes
éducatives complexes, l'orchestration devient essentielle. L'orchestration désigne la
logique de contrôle qui coordonne le flux d'information, délègue les sous-tâches, gère
l'état et garantit la cohérence des résultats finaux. Sans orchestration appropriée, les
agents risquent de boucler indéfiniment, d'entreprendre des actions contradictoires ou
d'utiliser les ressources de façon inefficace.

\subsection{Évolution des schémas d'orchestration}

L'orchestration d'agents a évolué rapidement depuis les premiers systèmes
d'enchaînement de prompts (\emph{prompt chaining}) en 2022--2023. Les systèmes initiaux reposaient sur
des pipelines linéaires. À partir de 2024, des schémas parallèles et hiérarchiques
sont apparus, poussés par la nécessité de fiabilité en production. En 2025--2026, les
environnements d'exécution basés sur des graphes sont devenus le paradigme dominant,
permettant des branchements conditionnels, un état persistant, de l'intervention
humaine dans la boucle et une replanification dynamique. Cette évolution répond
directement aux limites observées dans les premiers systèmes agentiques, notamment
dans les contextes éducatifs où les tâches impliquent de multiples services backend
et des données étudiantes en temps réel~\cite{anthropicagents,langgraphdocs}.

\subsection{Comparaison des principaux cadres d'agents}

Plusieurs frameworks dominent désormais l'écosystème. Le tableau ci-dessous résume
les différences clés pertinentes pour notre implémentation SprintUp :

\begin{figure}[H]
  \centering
  \includegraphics[width=1.0\textwidth]{generated-image.png}
  \caption{Comparaison des principaux frameworks d'agents IA (2025--2026)}
  \label{fig:frameworks-comparison}
\end{figure}

\subsection{LangGraph et environnement d'exécution basé sur des graphes}

\begin{figure}[H]
  \centering
  \includegraphics[width=0.32\textwidth]{logo-langraph.png}
  \caption{Logo officiel de LangGraph.}
  \label{fig:langgraph-logo}
\end{figure}

LangGraph (au sein de l'écosystème LangChain) illustre l'orchestration moderne. Il
représente les workflows sous forme de graphes orientés où les nœuds sont des agents
ou des outils et les arêtes définissent le routage conditionnel. Un état persistant et
des points de contrôle avec intervention humaine garantissent la fiabilité, un aspect
critique lors de l'intégration de données étudiantes sensibles~\cite{langgraphdocs}.
Dans notre couche SprintUp, nous tirerons parti du schéma de superviseur de LangGraph
combiné à l'exposition d'outils via MCP.

Trois principaux schémas d'orchestration dominent la pratique actuelle et sont bien
pris en charge par des cadres comme LangGraph. La figure~\ref{fig:sequential} présente
le premier schéma.

%% ── Schéma 1 : Séquentiel ───────────────────────────────────
\begin{figure}[H]
\centering
\begin{tikzpicture}[
    node/.style={rectangle, rounded corners=6pt, draw, thick,
      minimum width=2.7cm, minimum height=0.8cm, align=center},
    arrow/.style={->, thick, -{Stealth[length=3mm]}}
]
\node[node, fill=gray!10, font=\bfseries] (title) {Sequential};
\node[node, below=0.8cm of title]  (q)  {User Query};
\node[node, below=0.65cm of q]     (s1) {Step 1};
\node[node, below=0.65cm of s1]    (s2) {Step 2};
\node[node, below=0.65cm of s2]    (s3) {Step 3};
\node[node, below=0.65cm of s3]    (f)  {Final Answer};
\draw[arrow] (q) -- (s1);
\draw[arrow] (s1) -- (s2);
\draw[arrow] (s2) -- (s3);
\draw[arrow] (s3) -- (f);
\end{tikzpicture}
\caption{Schéma séquentiel d'orchestration}
\label{fig:sequential}
\end{figure}

Sequential orchestration organise les tâches dans une chaîne linéaire. Chaque étape
attend la fin de la précédente avant de transmettre sa sortie. Ce schéma convient bien
à des flux de travail clairement définis, comme \og{}récupérer l'emploi du temps du
cours $\rightarrow$ extraire les notes du chapitre concerné $\rightarrow$ générer des
exercices supplémentaires\fg{}. Sa simplicité facilite le débogage, mais il peut
devenir un goulot d'étranglement lorsque des sous-tâches pourraient s'exécuter
indépendamment.

%% ── Schéma 2 : Parallèle ────────────────────────────────────
La figure~\ref{fig:parallel} présente le deuxième schéma.

\begin{figure}[H]
\centering
\begin{tikzpicture}[
    node/.style={rectangle, rounded corners=6pt, draw, thick,
      minimum width=2.7cm, minimum height=0.8cm, align=center},
    arrow/.style={->, thick, -{Stealth[length=3mm]}}
]
\node[node, fill=gray!10, font=\bfseries]     (title) {Parallel};
\node[node, below=0.55cm of title]            (q)   {User Query};
\node[node, below=0.55cm of q]                (s)   {Splitter};
\node[node, below left=0.7cm and 0.3cm of s]  (a)   {Agent 1};
\node[node, below right=0.7cm and 0.3cm of s] (b)   {Agent 2};
\node[node, below=2.9cm of s]                 (m)   {Merge};
\node[node, below=0.55cm of m]                (f)   {Final Answer};
\draw[arrow] (q) -- (s);
\draw[arrow] (s) -- (a);
\draw[arrow] (s) -- (b);
\draw[arrow] (a) -- (m);
\draw[arrow] (b) -- (m);
\draw[arrow] (m) -- (f);
\end{tikzpicture}
\caption{Schéma parallèle d'orchestration}
\label{fig:parallel}
\end{figure}

Parallel (ou simultaneous) orchestration permet à plusieurs agents ou outils de
s'exécuter en même temps. Une requête peut lancer séparément la récupération de
l'emploi du temps du jour, la recherche de ressources multimédias et la rédaction
d'un résumé ; les résultats sont ensuite fusionnés par une étape de combinaison.
Cette approche améliore la rapidité et convient particulièrement aux contextes
éducatifs où les enseignants ont besoin d'un support rapide et multiple. La difficulté
principale réside dans la résolution des conflits et la synchronisation de l'état.

%% ── Schéma 3 : Supervisé ────────────────────────────────────
La figure~\ref{fig:supervised} présente le troisième schéma.

\begin{figure}[H]
\centering
\begin{tikzpicture}[
    node/.style={rectangle, rounded corners=6pt, draw, thick,
      minimum width=2.7cm, minimum height=0.8cm, align=center},
    arrow/.style={->, thick, -{Stealth[length=3mm]}}
]
\node[node, fill=gray!10, font=\bfseries] (title) {Supervised};
\node[node, below=0.8cm of title]                (q)   {User Query};
\node[node, below=0.9cm of q]                    (sup) {Supervisor};
\node[node, below left=1.05cm and -1.0cm of sup] (sp1) {Specialist 1};
\node[node, below right=1.05cm and -1.0cm of sup](sp2) {Specialist 2};
\node[node, below=2.2cm of sup]                  (f)   {Final Answer};
\draw[arrow] (q)   -- (sup);
\draw[arrow] (sup) -- (sp1);
\draw[arrow] (sup) -- (sp2);
\draw[arrow] (sp1) -- (sup);
\draw[arrow] (sp2) -- (sup);
\draw[arrow] (sup) -- (f);
\end{tikzpicture}
\caption{Schéma supervisé (hiérarchique) d'orchestration}
\label{fig:supervised}
\end{figure}

Supervised (ou hierarchical/orchestrator-worker) orchestration introduit un agent
superviseur central qui agit comme routeur et contrôleur de qualité. Le superviseur
analyse la requête utilisateur, décide quels agents spécialisés appeler, surveille
l'avancement, puis synthétise la réponse finale. Ce schéma, bien implémenté dans les
graphes superviseur de LangGraph, reflète la gestion d'une équipe humaine et offre
de bons mécanismes de gestion d'erreur et de repli. C'est le schéma que nous
privilégierons dans notre couche SprintUp, car il équilibre flexibilité et contrôle,
ce qui est essentiel pour une intégration avec des services backend en production.

LangGraph montre comment ces schémas peuvent coexister dans un même runtime basé sur
des graphes. Les nœuds représentent des agents ou des outils, les arêtes définissent
le flux de données et le routage conditionnel, et l'état persistant avec points de
contrôle humains assure la fiabilité. Ces cadres ont fait évoluer le développement
d'agents, passant de chaînes de prompts fragiles à des workflows testables et de
niveau production.

\subsection{Approfondissement du schéma de superviseur LangGraph}

Les environnements d'exécution basés sur des graphes comme LangGraph fournissent non
seulement des schémas abstraits, mais aussi des mécanismes concrets pour construire
des systèmes multi-agents robustes. Dans une configuration typique de superviseur, un
nœud contrôleur de haut niveau reçoit la requête utilisateur, inspecte l'état courant
et décide quels agents spécialistes ou outils invoquer
ensuite~\cite{langgraphsupervisor,langgraphdocs}. Ces spécialistes peuvent inclure par
exemple un agent \og{}cours\fg{}, un agent \og{}notes\fg{} ou un générateur
d'exercices, chacun responsable d'une capacité restreinte alignée sur un service backend
spécifique.

Le modèle de graphe d'état de LangGraph permet au superviseur de maintenir une mémoire
structurée des résultats intermédiaires, des sorties d'outils et des signaux d'erreur,
ce qui facilite la replanification lorsqu'un outil échoue ou que des informations
supplémentaires sont nécessaires~\cite{ibmlanggraph}. Les mécanismes de checkpointing et
les contrôles avec humain dans la boucle, désormais standards dans les versions modernes
de LangGraph, jouent un rôle clé dans les scénarios éducatifs où des interventions
d'enseignants ou d'administrateurs peuvent être requises avant l'exécution d'actions
sensibles sur les données étudiantes~\cite{anthropicagents}. Cette combinaison de
routage supervisé, d'état persistant et de points de contrôle explicites constituera
l'ossature de la couche d'orchestration SprintUp.

%% ────────────────────────────────────────────────────────────
\section{Le Model Context Protocol (MCP) : définition, origines et justification}

Tandis que l'orchestration gouverne le flux de haut niveau, le Model Context Protocol
traite la question plus bas niveau de la communication des agents avec les systèmes
externes. Annoncé par Anthropic en novembre 2024 puis étendu via des mises à jour de
spécification ouvertes, MCP est un protocole ouvert et standardisé qui fonctionne
comme un \og{}port USB-C de l'IA\fg{}~\cite{anthropic2024,anthropicmcpdocs}.

Il définit une architecture client--serveur cohérente. Côté client,
un composant MCP intégré dans l'environnement d'exécution de l'agent
charge les descriptions d'outils et leurs capacités dans la fenêtre
de contexte du modèle. Le modèle décide alors quels outils appeler
et émet des requêtes structurées. Côté serveur, des serveurs MCP
exposent les services backend (bases de données, API, systèmes de
fichiers, fonctions métier) de façon uniforme via JSON-RPC sur des
canaux sécurisés. La spécification 2025 définit deux transports
principaux : \texttt{stdio} pour les intégrations locales et
\textbf{Streamable HTTP} (successeur du transport SSE initial) pour
les déploiements réseau ; ce dernier permet des requêtes HTTP POST
stateless avec réponses JSON directes, facilitant le passage à
l'échelle derrière un load balancer. Les résultats transitent par
le même protocole, ce qui autorise des interactions multi-tours
fluides.

À la différence des mécanismes traditionnels d'appel de fonctions
(\emph{function calling}) qui exigent du code de liaison spécifique pour
chaque nouveau service, MCP découple l'agent des
détails d'implémentation. Les développeurs implémentent le serveur MCP une seule fois
pour leur backend (dans notre cas, les endpoints de services étudiants pour les cours,
notes, génération de contenu, etc.), et tout agent conforme peut ensuite découvrir et
utiliser immédiatement ces outils. MCP a précisément été conçu pour résoudre la
fragmentation qui pénalisait les premiers développements d'agents : avant 2024, chaque
entreprise devait écrire des intégrations sur mesure pour chaque base de données, API
ou outil, un cauchemar de maintenance qui passait mal à l'échelle et introduisait des
incohérences de sécurité~\cite{anthropic2024,mcpspec}. D'un point de vue
cybersécurité, au cœur de ce programme, les hooks d'authentification intégrés au
protocole, les périmètres d'autorisation et le format de messages propice à l'audit
représentent une avancée majeure par rapport aux solutions \emph{ad hoc}
d'appel de fonctions~\cite{mcpsecurity}.

MCP ne remplace pas la récupération d'information ; il la complète. La recherche
agentique sur sources externes excelle dans la collecte de connaissances
actualisées pour ancrer le contenu généré dans des sources fiables. MCP permet à
l'agent d'agir sur ces connaissances en appelant des services applicatifs en temps réel.
Dans le contexte SprintUp, la rédaction d'une activité déclenche d'abord la
collecte du contexte pédagogique pertinent par recherche externe, puis les
outils \gls{mcp} de création et de mise à jour pour persister le contenu dans
la base de connaissances.

\subsection{Architecture détaillée et flux protocolaire}

MCP définit une architecture client--serveur basée sur JSON-RPC via des événements
envoyés par le serveur. Le client MCP (intégré dans l'exécution de l'agent) charge les
descriptions d'outils dans la fenêtre de contexte du modèle et les expose sous forme
de définitions structurées que le modèle peut appeler~\cite{anthropicmcpdocs,mcpspec}.
Quand le modèle décide d'invoquer un outil, il émet une requête JSON-RPC :

\begin{tcolorbox}[colback=gray!5, colframe=gray!50, arc=3pt,
    title={\small Exemple de requ\^{e}te JSON-RPC MCP}]
\small\ttfamily
\{``jsonrpc'': ``2.0'', ``method'': ``tools/call'',\\
\ ``params'': \{``name'': ``create\_activity'',\\
\quad``arguments'': \{``unit\_id'': ``u-01'',\\
\qquad``title'': ``Les bases de données relationnelles''\}\}\}
\end{tcolorbox}

Le serveur MCP reçoit la requête, exécute la logique backend correspondante
et renvoie une réponse structurée :

\begin{tcolorbox}[colback=gray!5, colframe=gray!50, arc=3pt,
    title={\small Exemple de r\'{e}ponse JSON-RPC MCP}]
\small\ttfamily
\{``jsonrpc'': ``2.0'', ``id'': 1,\\
\ ``result'': \{``id'': ``a-42'',\\
\quad``title'': ``Les bases de données relationnelles'',\\
\quad``position'': 0, ``status'': ``created''\}\}
\end{tcolorbox}

Le transport \textbf{\emph{Streamable HTTP}} se prête naturellement à un mode
sans état dans lequel chaque appel \gls{rpc} est auto-contenu et n'oblige pas
le serveur à mémoriser une session de
client~\cite{anthropicmcpdocs}. La figure~\ref{fig:mcp-architecture} illustre
cette architecture de haut niveau.

\begin{figure}[H]
\centering
\begin{tikzpicture}[
    node/.style={rectangle, rounded corners=6pt, draw, thick,
      minimum width=3.5cm, minimum height=1cm, align=center},
    arrow/.style={->, thick, -{Stealth[length=3mm]}}
]
\node[node] (model)  {Agent LLM};
\node[node, below=1.5cm of model]  (client) {Client MCP};
\node[node, below=1.5cm of client] (server) {Serveur MCP (Backend)};
\draw[arrow] (model)  -- node[midway,right] {Appel d'outil} (client);
\draw[arrow] (client) -- node[midway,right] {JSON-RPC}      (server);
\draw[arrow] (server) -- node[midway,left]  {Réponse}       (client);
\draw[arrow] (client) -- node[midway,left]  {Résultat}      (model);
\end{tikzpicture}
\caption{Architecture client--serveur MCP (spécification 2025)}
\label{fig:mcp-architecture}
\end{figure}

Du point de vue de l'implémentation, la spécification MCP et les implémentations de
référence dans des langages courants (Python, TypeScript) fournissent une cartographie
claire entre les capacités côté serveur et les schémas d'outils visibles dans le
contexte du modèle~\cite{anthropicmcpdocs}. Cela réduit la complexité accidentelle et
facilite la compréhension des opérations qu'un agent est autorisé à effectuer dans un
déploiement éducatif donné.

\subsection{Fonctionnalités de sécurité et risques}

Sur le plan de la cybersécurité, MCP présente plusieurs avantages par rapport aux
schémas informels d'appel de fonctions. Tout d'abord, il encourage une définition
explicite des outils disponibles et de leurs paramètres, ce qui s'aligne sur le
principe du moindre privilège : un agent interagissant avec SprintUp peut n'avoir
accès qu'à des méthodes de consultation de cours en lecture seule pour les profils
étudiants, tandis que les agents enseignants bénéficient de capacités d'écriture plus
larges~\cite{mcpsecurity}. Ensuite, les messages JSON-RPC structurés constituent une
piste d'audit naturelle pouvant être journalisée et surveillée pour détecter des
comportements suspects.

Des risques subsistent néanmoins. Des attaquants pourraient tenter d'enregistrer des
serveurs MCP malveillants exposant des outils conçus pour exfiltrer des données ou
injecter du contenu adversarial dans le contexte de l'agent. Des attaques par injection
de prompt ou empoisonnement de données peuvent également survenir via les réponses
d'outils si celles-ci ne sont pas correctement filtrées~\cite{mcpsecurity}.
Plusieurs mesures de mitigation sont recommandées. La première
consiste à exiger une authentification et une autorisation strictes
des serveurs MCP avant leur mise à disposition de l'agent. La
deuxième repose sur des périmètres fins et un contrôle d'accès basé
sur les rôles (\acrshort{rbac}) définissant les outils qu'un profil
d'utilisateur peut invoquer. La troisième impose la validation des
entrées et des sorties de tous les appels d'outils, avec filtrage
du contenu non fiable avant qu'il n'atteigne le modèle. La
quatrième, enfin, repose sur une journalisation exhaustive et une
détection d'anomalies sur le trafic MCP, en particulier pour les
opérations impliquant des données personnelles d'étudiants.

Dans l'implémentation SprintUp, ces mesures seront intégrées aux contrôles de sécurité
existants et aux politiques de protection des données.

%% ────────────────────────────────────────────────────────────
\section{État de l'art : applications en éducation}

La combinaison de l'orchestration, de la Retrieval-Augmented Generation (RAG) et
d'intégrations de type MCP a déjà donné lieu à des résultats convaincants dans des
environnements éducatifs à travers le monde. L'Arizona State University (ASU), par
exemple, a lancé en 2024 un important \og{}AI Innovation Challenge\fg{} ayant attiré
plus de 175 propositions et financé plus de 200 projets couvrant la majorité de ses
unités académiques~\cite{asu2024}. Ces initiatives incluent des agents de tutorat
personnalisé et des outils de simulation combinant RAG pour la récupération de
connaissances et des mécanismes d'orchestration pour fournir des expériences
d'apprentissage adaptatives.

De la même façon, des prototypes basés sur RAG ont montré de bonnes performances dans
les environnements universitaires. Une revue systématique récente sur RAG pour les
applications éducatives met en évidence la façon dont les modèles augmentés par la
récupération peuvent alimenter des systèmes d'apprentissage interactifs, la génération
de contenu et l'évaluation automatisée, tout en réduisant les hallucinations par rapport
aux modèles purement génératifs~\cite{rageducation,lirageducation}.

Un exemple particulièrement pertinent pour ce projet est l'expérience d'Intellect
Academy en Tunisie, qui opère via la plateforme SprintUp sur plusieurs campus et
filières de formation~\cite{intellect2025}. En intégrant l'IA pour le développement de
curricula, les parcours d'apprentissage adaptatifs et le pilotage unifié de la
performance, l'académie a accéléré la création de contenu tout en préservant une forte
maîtrise pédagogique.

L'Université Northeastern a, elle, noué un partenariat avec Anthropic pour déployer des
agents basés sur Claude sur ses campus pour le tutorat personnalisé et le support
administratif~\cite{northeastern2025,duke2025}. Ces initiatives illustrent une tendance
plus large en 2025--2026 : l'émergence d'universités \og{}IA natives\fg{} où des
systèmes agentiques soutiennent le conseil, le tutorat et les workflows d'alerte
précoce à grande échelle~\cite{agenticragsurvey,agenticrag2025}.

\subsection{RAG agentique pour l'apprentissage personnalisé}

Les pipelines classiques de Retrieval-Augmented Generation (RAG) récupèrent un ensemble
statique de documents à partir de la requête utilisateur, puis génèrent une réponse en
un seul passage. Si cela améliore l'ancrage factuel, ces workflows peuvent rencontrer
des difficultés pour le raisonnement multi-étapes, les questions ambiguës ou les tâches
nécessitant une exploration adaptative de la base de
connaissances~\cite{lewis2020,rageducation}. Le RAG agentique étend ce paradigme en
intégrant des agents autonomes dans la chaîne RAG. Ces agents s'appuient sur des
schémas de conception tels que la planification, la réflexion, l'utilisation d'outils
et la collaboration multi-agents pour piloter dynamiquement les stratégies de
récupération, affiner itérativement le contexte et adapter les
workflows~\cite{agenticragsurvey}.

Dans des scénarios d'apprentissage personnalisé, un système de RAG
agentique peut ainsi planifier une séquence de récupérations à
partir de notes de cours, de chapitres de manuels et d'anciens
sujets d'examen, puis réfléchir à la suffisance du contexte
récupéré et émettre des requêtes complémentaires en cas de lacunes.
Il peut également utiliser des outils pour parcourir des graphes
de connaissances ou interroger un modèle de l'étudiant afin
d'adapter les explications à son niveau actuel.

La revue consacrée à RAG pour les applications éducatives met spécifiquement en avant
quatre grands schémas d'application : systèmes de questions--réponses éducatifs,
chatbots pédagogiques, systèmes de tutorat pilotés par IA et parcours d'apprentissage
adaptatifs~\cite{rageducation,lirageducation}. Chacun de ces schémas s'aligne
étroitement avec les services backend existants de SprintUp.

%% ────────────────────────────────────────────────────────────
\section{Conclusion du chapitre}

Ce chapitre a posé les bases conceptuelles : l'orchestration fournit l'intelligence de
coordination des tâches complexes, tandis que MCP apporte la plomberie standardisée et
sécurisée qui connecte les agents aux systèmes backend réels. Combinées à RAG (détaillé
dans les chapitres suivants), ces technologies constituent la pile complète nécessaire
pour concrétiser notre vision pour SprintUp. Les schémas et les protocoles présentés ici
ne sont pas théoriques ; ils adressent directement le déficit d'intégration mis en
évidence en introduction et orienteront chacune des décisions de conception à venir.

\subsection{Synthèse et lien avec SprintUp}

Le schéma d'orchestration supervisée, combiné à l'exposition d'outils via MCP et à
l'ancrage par RAG, s'ajuste naturellement aux services backend de SprintUp. L'analyse de
la sécurité de MCP et de l'orchestration d'agents présentée dans ce chapitre fournit un
socle pour concevoir une couche d'agents conformes et auditables, respectant la vie
privée des étudiants et les contraintes institutionnelles. Forts de ces fondations, nous
pouvons désormais passer à l'implémentation concrète de la récupération d'information,
de l'exposition des outils via MCP et de la logique d'orchestration dans les chapitres
de mise en \oe{}uvre qui suivent.
\newpage
